\begin{proof}

\begin{enumerate}[label={$\langle1\rangle$\arabic*.},leftmargin=*]
\item Each nontrivial degree-zero line on an elliptic curve has no
sections: a nonzero section would have an effective divisor of
degree zero and hence trivialize it. Riemann--Roch and Serre
duality then give zero first cohomology. The trivial line has
cohomology $(1,1)$. To compute the direct image with its length,
put $p:A\times A\to A$ and $M=P\otimes\operatorname{pr}_1^*
\mathcal O_A(0)$. Its fibres have degree one, so
$R^1p_*M=0$ and $p_*M$ is a line bundle. The exact sequence for
the section $\{0\}\times A$ gives the two-term complex
\[
 Rp_*P=\bigl[p_*M\longrightarrow\mathcal O_A\bigr].
\]
The arrow evaluates the unique section of the moving degree-one
line at the origin. Its zero is exactly $t=0$ and is simple:
in a local coordinate the moving zero divisor is the diagonal,
with equation $x-t$. Thus the arrow is the inclusion
$\mathcal O_A(-0)\to\mathcal O_A$, and
$Rp_*P=\mathbb C_0[-1]$. Tensoring with
$R\Gamma(\mathbb P^1,\mathcal O(-2))=\mathbb C[-1]$
gives $Rg_*V=\mathbb C_0[-2]$. This proves the direct-image,
length, and global-cohomology assertions. Fibre Kunneth gives
the stated fibre cohomology.

\item For any proper flat family over a nonsingular curve and a locally
free sheaf $V$, multiplication by a local parameter gives the
actual replacement exact sequence
\[
0\to (R^qg_*V)_c/t(R^qg_*V)_c
 \to H^q(Y_c,V|_{Y_c})
 \to (R^{q+1}g_*V)_c[t]\to0.
\]
It follows directly from the long direct-image sequence of
$0\to V\xrightarrow{t}V\to V|_{Y_c}\to0$. Here $[t]$ denotes
the kernel of multiplication by $t$. Local constancy of all
fibre cohomology dimensions implies local freeness and base
change by Grauert's theorem; generic vanishing alone permits
the torsion explicitly present above.

\end{enumerate}
\end{proof}
